\subsection{Comparison with State-of-the-art Methods}
    \label{subsec:comparisons}
We compare GDKVM with four task-specific methods (PKEchoNet~\cite{wu2023super}, DSA~\cite{lin2024dynamic}, MemSAM~\cite{deng2024memsam}, SimLVSeg~\cite{maani2024simlvseg}) and four related approaches (Xmem++~\cite{bekuzarov2023xmemproductionlevelvideosegmentation}, Cutie~\cite{cheng2024puttingobjectvideoobject}, VideoMamba~\cite{li2024videomambastatespacemodel}, Vision LSTM~\cite{beck:24xlstm}). 

These prior works incorporate memory mechanisms, time-series modeling, or efficient attention modules and have advanced medical video segmentation. However, they still face challenges in capturing complex spatial variations and retrieving temporal information accurately. GDKVM addresses these limitations by introducing key-value associations within a linear matching framework, which reduces the typical high computational overhead of conventional attention. In \cref{tab:cap}, GDKVM delivers higher overall metrics on two datasets and outperforms four widely used task-specific models. In \cref{fig:fps}, GDKVM achieves an mDice of 95.11 at 37 FPS with 35.2~M parameters. GDKVM leverages GDR for dynamic memory management to discard irrelevant information promptly and update essential memory content efficiently, matching or surpassing specialized models in segmentation performance. This design avoids bottlenecks arising from incomplete memory mechanisms or overly complex training procedures, allowing GDKVM to maintain both accuracy and speed.